% GENERATED FILE. Edit canonical lessons, then run npm run generate:books.
% canonical-source-hash: fnv1a64:966c3ef4dc995447
% canonical-lessons: IT-C278-principio, IT-C278-terminare, IT-C278-concludere, IT-C278-conclusione, IT-C278-riprendere

\chapter{Dal principio alla fine}
\label{ch:it-dal-principio-alla-fine}
\hlchaptermodality{278}

\begin{chapteropening}
\textbf{By the end of this chapter:} \emph{I can say how something began, started again, and came to an end.}

Complete the last lesson of chapter 278: I can say how something began, started again, and came to an end.
\end{chapteropening}

\section[il principio]{il principio — the beginning; the principle}

\label{lesson:IT-C278-principio}



\begin{quote}
Before the new one: say the Italian for the moment, then the Italian for the phase.
\end{quote}

\subsection*{You'll want to know: il principio}
\textbf{il principio} — \textquotedblleft{}the beginning; the principle\textquotedblright{}.

\begin{quote}
Raccontami tutto dal principio.
\end{quote}

\begin{quote}
\emph{Tell me everything from the beginning.}
\end{quote}

\emph{In principio} — in the beginning. The plural \emph{i principi} means principles.

\begin{cousinweb}
\emph{principio} is a learned borrowing from Latin \emph{prīncipium} \textquotedblleft{}a beginning, a foundation\textquotedblright{} (EtymDB), built on \emph{prīnceps} \textquotedblleft{}first\textquotedblright{}.
\end{cousinweb}

\begin{grammarlens}[title={What you've built}]
One more word for telling, explaining and describing at B1.
\end{grammarlens}

\subsection*{Your turn}
\begin{itemize}
\raggedright
  \item \emph{Say it:} \emph{il principio}
  \item \emph{Say it:} \emph{il principio}, once more
  \item \emph{Say it:} a sentence of your own with \emph{il principio}
  \item \emph{Recall:} say the Italian for the moment, then the Italian for the phase, then say \emph{il principio} again
\end{itemize}

\begin{tcolorbox}[breakable,colback=teal!4,colframe=teal!35!black,title={Before you move on}]
What does il principio mean? (\textquotedblleft{}The beginning; the principle\textquotedblright{}.) Say it once more.
\end{tcolorbox}
\section[terminare]{terminare — to finish, to end}

\label{lesson:IT-C278-terminare}



\begin{quote}
Before the new one: say the Italian for the phase, then the Italian for the beginning.
\end{quote}

\subsection*{You'll want to know: terminare}
\textbf{terminare} — \textquotedblleft{}to finish, to end\textquotedblright{}.

\begin{quote}
La lezione termina alle sei.
\end{quote}

\begin{quote}
\emph{The lesson ends at six.}
\end{quote}

A formal \emph{finire}; with or without an object: \emph{ho terminato il lavoro}, \emph{lo spettacolo è terminato}.

\begin{cousinweb}
\emph{terminare} is a learned borrowing from Latin \emph{termināre} \textquotedblleft{}to mark the bounds; to end\textquotedblright{} (EtymDB), built on \emph{terminus} \textquotedblleft{}boundary, limit\textquotedblright{}.
\end{cousinweb}

\begin{grammarlens}[title={What you've built}]
One more word for telling, explaining and describing at B1.
\end{grammarlens}

\subsection*{Your turn}
\begin{itemize}
\raggedright
  \item \emph{Say it:} \emph{terminare}
  \item \emph{Say it:} \emph{terminare}, once more
  \item \emph{Say it:} a sentence of your own with \emph{terminare}
  \item \emph{Recall:} say the Italian for the phase, then the Italian for the beginning, then say \emph{terminare} again
\end{itemize}

\begin{tcolorbox}[breakable,colback=teal!4,colframe=teal!35!black,title={Before you move on}]
What does terminare mean? (\textquotedblleft{}To finish, to end\textquotedblright{}.) Say it once more.
\end{tcolorbox}
\section[concludere]{concludere — to conclude, to wind up}

\label{lesson:IT-C278-concludere}



\begin{quote}
Before the new one: say the Italian for the beginning, then the Italian for to finish.
\end{quote}

\subsection*{You'll want to know: concludere}
\textbf{concludere} — \textquotedblleft{}to conclude, to wind up\textquotedblright{}.

\begin{quote}
Per concludere, vi ringrazio tutti.
\end{quote}

\begin{quote}
\emph{To conclude, I thank you all.}
\end{quote}

Past participle \emph{concluso}: \emph{abbiamo concluso l'accordo}.

\begin{cousinweb}
\emph{concludere} is a learned borrowing from Latin \emph{conclūdere} \textquotedblleft{}to shut in; to close, to end\textquotedblright{} (EtymDB), from \emph{con-} and \emph{claudere} \textquotedblleft{}to shut\textquotedblright{}.
\end{cousinweb}

\begin{grammarlens}[title={What you've built}]
One more word for telling, explaining and describing at B1.
\end{grammarlens}

\subsection*{Your turn}
\begin{itemize}
\raggedright
  \item \emph{Say it:} \emph{concludere}
  \item \emph{Say it:} \emph{concludere}, once more
  \item \emph{Say it:} a sentence of your own with \emph{concludere}
  \item \emph{Recall:} say the Italian for the beginning, then the Italian for to finish, then say \emph{concludere} again
\end{itemize}

\begin{tcolorbox}[breakable,colback=teal!4,colframe=teal!35!black,title={Before you move on}]
What does concludere mean? (\textquotedblleft{}To conclude, to wind up\textquotedblright{}.) Say it once more.
\end{tcolorbox}
\section[la conclusione]{la conclusione — the conclusion}

\label{lesson:IT-C278-conclusione}



\begin{quote}
Before the new one: say the Italian for to finish, then the Italian for to conclude.
\end{quote}

\subsection*{You'll want to know: la conclusione}
\textbf{la conclusione} — \textquotedblleft{}the conclusion\textquotedblright{}.

\begin{quote}
Alla fine siamo arrivati a una conclusione.
\end{quote}

\begin{quote}
\emph{In the end we reached a conclusion.}
\end{quote}

\emph{In conclusione} — in conclusion.

\begin{cousinweb}
\emph{conclusione} is a learned borrowing from Latin \emph{conclūsiō} \textquotedblleft{}a closing, a conclusion\textquotedblright{} (EtymDB), from \emph{conclūdere}.
\end{cousinweb}

\begin{grammarlens}[title={What you've built}]
One more word for telling, explaining and describing at B1.
\end{grammarlens}

\subsection*{Your turn}
\begin{itemize}
\raggedright
  \item \emph{Say it:} \emph{la conclusione}
  \item \emph{Say it:} \emph{la conclusione}, once more
  \item \emph{Say it:} a sentence of your own with \emph{la conclusione}
  \item \emph{Recall:} say the Italian for to finish, then the Italian for to conclude, then say \emph{la conclusione} again
\end{itemize}

\begin{tcolorbox}[breakable,colback=teal!4,colframe=teal!35!black,title={Before you move on}]
What does la conclusione mean? (\textquotedblleft{}The conclusion\textquotedblright{}.) Say it once more.
\end{tcolorbox}
\section[riprendere]{riprendere — to resume, to start again; to take back}

\label{lesson:IT-C278-riprendere}



\begin{quote}
Before the new one: say the Italian for to conclude, then the Italian for the conclusion.
\end{quote}

\subsection*{You'll want to know: riprendere}
\textbf{riprendere} — \textquotedblleft{}to resume, to start again; to take back\textquotedblright{}.

\begin{quote}
Dopo la pausa, abbiamo ripreso a lavorare.
\end{quote}

\begin{quote}
\emph{After the break, we started working again.}
\end{quote}

Past participle \emph{ripreso}. \emph{Riprendere a} + infinitive: to start doing again.

\begin{cousinweb}
\emph{riprendere} is inherited from Latin \emph{reprehendere} \textquotedblleft{}to hold back, to seize again\textquotedblright{} (EtymDB), built on \emph{prehendere} \textquotedblleft{}to take hold of\textquotedblright{}.
\end{cousinweb}

\begin{grammarlens}[title={What you've built}]
One more word for telling, explaining and describing at B1.
\end{grammarlens}

\subsection*{Your turn}
\begin{itemize}
\raggedright
  \item \emph{Say it:} \emph{riprendere}
  \item \emph{Say it:} \emph{riprendere}, once more
  \item \emph{Say it:} a sentence of your own with \emph{riprendere}
  \item \emph{Recall:} say the Italian for to conclude, then the Italian for the conclusion, then say \emph{riprendere} again
\end{itemize}

\begin{tcolorbox}[breakable,colback=teal!4,colframe=teal!35!black,title={Before you move on}]
What does riprendere mean? (\textquotedblleft{}To resume, to start again; to take back\textquotedblright{}.) Say it once more.
\end{tcolorbox}
